\begin{tabular}{lrll}
\toprule
параметр & значение & интерпретация & источник или комментарий \\
\midrule
beta & 0.972633 & Коэффициент дисконтирования & Подбирается по asset market \\
eis & 0.500000 & Межвременная эластичность замещения & Baseline two-asset HANK \\
frisch & 1.000000 & Фришева эластичность труда & Baseline two-asset HANK \\
chi0 & 0.350000 & Сдвиг в знаменателе функции издержек ребалансировки & Baseline two-asset HANK \\
chi1 & 23.000000 & Масштаб издержек ребалансировки & Подбирается в steady state \\
chi2 & 2.000000 & Кривизна издержек ребалансировки & Baseline two-asset HANK \\
phi_pi & 1.500000 & Коэффициент реакции правила Тейлора на инфляцию & Классическая денежно-кредитная политика \\
phi_y & 0.125000 & Коэффициент реакции правила Тейлора на выпуск & Классическая денежно-кредитная политика \\
rho_i & 0.700000 & Инерционность процентной ставки & Классическая денежно-кредитная политика \\
kappap & 0.100000 & Наклон ценовой кривой Филлипса & Номинальные жесткости \\
kappaw & 0.100000 & Наклон кривой Филлипса по заработной плате & Номинальные жесткости \\
delta & 0.020000 & Норма амортизации капитала & Производственный блок \\
omega & 0.015000 & Спрэд между доходностью ликвидного актива и ставкой политики & Клин ликвидной доходности \\
rho_z & 0.966000 & Персистентность идосинкратического дохода & Процесс идосинкратического дохода \\
sigma_z & 0.880000 & Безусловная дисперсия лог-идосинкратического дохода & Процесс идосинкратического дохода \\
nB & 18.000000 & Число узлов по liquid asset & Численная сетка \\
nA & 28.000000 & Число узлов по illiquid asset & Численная сетка \\
nZ & 3.000000 & Число узлов доходного процесса & Численная сетка \\
shock & 0.001000 & Размер монетарного шока & Эксперимент денежно-кредитной политики \\
rho_eps & 0.000000 & Персистентность монетарного шока & Эксперимент денежно-кредитной политики \\
\bottomrule
\end{tabular}
